\properlane{word-received-in-trial}{The word received in trial becomes a witness to others}

\subsection{What the disciples have just heard}

The promise of growth follows some discouraging teaching. In the parable of the sower, several kinds of ground fail to produce a harvest. In the parable of the weeds, the crop that does grow suffers the enemy's interference. Chrysostom imagines the question these accounts can raise in the disciples: how many faithful people will there be? He reads the mustard seed as Christ's answer to that fear. Their present weakness does not determine the reach of the preaching entrusted to them.\footnote{Chrysostom, \textit{Homilies on Matthew}, 46.2.}

That explanation depends on the sequence of Matthew 13. The question is not simply how something small becomes large. It is how the disciples can trust the work they are undertaking after hearing Christ describe failure and opposition. The seed's growth answers discouragement without requiring them to forget those dangers. A plant with lodging for birds is a larger result than its beginning suggests; its development gives the disciples an image through which to understand Christ's power at work in their weakness.

Chrysostom then reads the leaven as a second explanation of the same assurance. Its small quantity is placed within much meal. The woman hides it there, and that contact is necessary to its action. Leaven kept apart from flour changes no dough. Chrysostom applies this detail directly to the disciples among the multitude. Their contact with opponents exposes them to injury, but it also makes their teaching effective among people who have not yet received it. The image permits him to address the fear that closeness itself will destroy the messengers.

His interpretation of ``three measures'' serves that argument. He takes the expression to indicate a large quantity. The force of the comparison lies in the disproportion between the little leaven and the much greater amount it affects. He does not assign each measure a separate person, faculty or stage of spiritual life. This is why his explanation can sustain a communal interpretation without requiring an allegorical identification for every item in the kitchen. The material action already explains the missionary point: Christ's power works through disciples who remain among those to whom they are sent.\footnote{Chrysostom, \textit{Homilies on Matthew}, 46.2, on Matthew 13:33.}

The human situation calls for judgment as well as courage. Chrysostom's leaven communicates what it has received; its effectiveness consists in changing the meal into its own condition. The application to Christian witness therefore concerns faith communicated through life and teaching. It gives no reason to call every form of influence evangelical. A community may become conspicuous through wealth or resentment without showing the virtues Paul praises. The Epistle makes the character of the desired fruit more definite.

\subsection{The difference between attracting attention and becoming an example}

Aquinas explains the Thessalonian thanksgiving by beginning with its recipient: thanks are addressed to God, are continued, and embrace the whole community. Paul both recognizes the believers' good conduct and asks for their further good. Thanksgiving and intercession can concern the same people because grace already received does not make further grace unnecessary.\footnote{Aquinas, \textit{Commentary on First Thessalonians}, chapter 1, lecture 1 [87862], on 1:2--3.}

His account of faith is deliberately active. He joins work and labor to faith because assent must issue in action and endure when action becomes difficult. He then speaks of the Thessalonians' abundant charitable works, and of hope as the reason they patiently bear adversity. Each virtue supplies a different explanation of what Paul sees. Faith directs action toward what has been believed; charity accounts for service to others; hope permits a person to remain faithful when the promised good has not yet appeared. This develops the Collect's request that attention be carried into words and deeds. In this interpretation, the prayer asks for the integrity through which a community's confession becomes recognizable in its conduct.

The Epistle also credits the Gospel's arrival to power, the Holy Spirit and fullness. Aquinas offers two explanations. The words may concern the confirmation of preaching through miracles and the gift of the Spirit. They may also concern the way Paul preached: a virtuous life, speech prompted by the Spirit, and instruction adequate to the faith. The alternatives enlarge the account of divine action. Neither permits the preacher to identify that action simply with his own persuasive ability.\footnote{Aquinas, \textit{Commentary on First Thessalonians}, chapter 1, lecture 1 [87862], on 1:4--5.}

This distinction matters when Chrysostom turns from the apostles to his own congregation. Some listeners could excuse themselves by saying that the apostles performed miracles. He answers by examining the life of the apostles: freedom from attachment to wealth, refusal of vainglory and mastery of the passions. Signs alone do not make the worker holy. A person can possess striking gifts and still act unjustly. Chrysostom consequently asks his hearers to practice virtues within their own reach rather than admire extraordinary powers as a reason for doing nothing.\footnote{Chrysostom, \textit{Homilies on Matthew}, 46.3.}

Aquinas supplies a further safeguard when he explains imitation. Christians imitate their teachers insofar as those teachers imitate Christ. Human failings do not acquire authority merely because they occur in a leader. In the appointed passage, the specific imitation concerns reception of the word amid suffering. Those who first followed the apostolic example become examples to other believers. Their conversion is reported before Paul needs to describe it. Their life has furnished other people with something concrete to recount.\footnote{Aquinas, \textit{Commentary on First Thessalonians}, chapter 1, lecture 1 [87862], on 1:6--9.}

That is the connection proposed here between the two lessons. Chrysostom explains why a few disciples need not withdraw in fear; Aquinas explains what fruitful reception looks like in a community. The Thessalonians have been changed in ways others can see. Their example also helps those who already believe, since Paul addresses its effect on believers in neighboring regions. Evangelical fruit includes strengthening existing faith as well as calling people to conversion.

\subsection{Joy while the affliction remains}

The unusual phrase in Paul's description is joy in much tribulation. Chrysostom pauses over the apparent difficulty: how can affliction and joy occur together? His answer distinguishes the harm suffered from what the Spirit brings about within those who suffer for Christ. Affliction is painful; the joy comes from the Holy Spirit's presence and their relation to Christ. He compares their condition to the biblical children refreshed within the furnace. The fire's ordinary effect does not explain their relief.\footnote{Chrysostom, \textit{Homilies on First Thessalonians}, homily 1, on 1:6.}

This is a claim about the Spirit's action, not a test requiring a distressed person to exhibit a cheerful manner. The Thessalonians' suffering remains named as suffering. Their fidelity shows that it has not separated them from the Lord whom they received. Aquinas relates their joy to love: because they love Christ, they can bear hardship for his sake. That explanation permits sorrow at what has happened and gratitude for belonging to Christ to exist together. It also explains why the example can encourage others facing similar difficulties.

The Introit's promise of peace can be heard alongside this experience. Jeremiah addresses people whose restoration still requires waiting. Jerome's exposition makes their prayer and whole-hearted seeking part of that waiting. A community formed by the Gospel can therefore ask for deliverance without making immediate relief the condition of its fidelity. The Gradual's promise of enduring praise likewise need not describe a congregation already beyond danger. Augustine's interpretation of its future rescue gives the praise an end that present opposition cannot determine.\footnote{Jerome, \textit{Commentary on Jeremiah}, V, on 29:10--14, PL 24, cols.~893--894; Augustine, \textit{Expositions of the Psalms}, 43.6--8.}

Schuster hears the Gradual in relation to martyrs whose faith has survived what destroys their bodily life. His application names a kind of victory compatible with the experience of persecution.\footnote{Schuster, \textit{The Sacramentary}, III, p.~187, commentary on \textit{Liberasti nos}.} It does not require calling every human loss a success. The relevant good is perseverance in the faith, and its promised completion remains with God. This also preserves the full psalm's serious complaint: the faithful may describe what hurts and ask God to act.

\subsection{The community continues to ask}

The Alleluia and Offertory give that asking a repeated expression. The community that has received the word still calls from the depths. In Augustine's exposition, prayer can reach God even where bodily escape is impossible. The petition thus suits those whose circumstances remain constrained. Its repetition at the Offertory places dependence within the act of offering. Schuster's explanation of gifts first received from God helps describe the assembly's action: it brings what his generosity has enabled it to bring, and it asks him to receive it.\footnote{Augustine, \textit{Expositions of the Psalms}, 129.1; Schuster, \textit{The Sacramentary}, III, pp.~187--188, on the Alleluia and Offertory.}

The Secret then asks for changes in the people offering. Cleansing matters because an unrepented wrong can contradict the truth they profess. Renewal matters because sustained witness requires a life continually received from God. Governance asks that their efforts be rightly directed; protection asks that the good already given endure. These are implications of the four petitions for a community whose conduct affects others. Schuster identifies the petitions as effects of the sacrifice; the connection with exemplary witness follows from reading them with Paul's thanksgiving.

Mission also has an object beyond the enlargement of a group. In Paul's account, conversion ends in serving the living and true God. The Trinity Preface is an act of that worship: the Father is thanked with the Son and Holy Spirit confessed as one God. The angels' praise enlarges the company with whom the assembly gives thanks. Aquinas's exposition of the letter's salutation relates its naming of Father and Son to faith in the Trinity and to the blessedness sought in the knowledge of God.\footnote{Aquinas, \textit{Commentary on First Thessalonians}, chapter 1, lecture 1 [87862], on 1:1; \textit{Missale Romanum} (1962), p.~293.} The salutation precedes the appointed thanksgiving, but it explains the identity of the community Paul addresses.

The Communion encourages confidence in prayer. Mark's accompanying command to forgive makes that confidence pertinent to relations within and beyond the community. Bede explains faith through charity and tells those seeking pardon to forgive others. Witness to the Gospel would become incoherent if those asking mercy refused it to their neighbor. Prayer for help in a difficult work includes consent to the moral obligations Christ places on the person praying.\footnote{Bede, \textit{Commentary on Mark}, on Mark 11:22--25, especially \textit{Et cum stabitis ad orandum}.}

Finally, the Postcommunion asks that nourishment generate an enduring appetite for the gifts received. The assembly's life will continue beyond this celebration, and its work, speech and endurance will require continued dependence on God. Its prayer does not assume that one fruitful moment has made such dependence unnecessary. Desire for the heavenly food belongs to the perseverance that Paul praises.

\subsection{The future for which the Church waits}

Paul ends with the Son from heaven. Aquinas identifies two objects of expectation: resurrection, in conformity with the risen Christ, and deliverance from the punishment that threatens sinners. He relates the latter to Christ's freeing them from sin.\footnote{Aquinas, \textit{Commentary on First Thessalonians}, chapter 1, lecture 1 [87862], on 1:10.} The community's influence in neighboring regions is therefore a present fruit, while its final hope concerns life with the risen Lord. An account of growth that ended with reputation would stop before Paul's sentence ends.

The Gospel's ``whole'' likewise belongs within Matthew's complete discourse. That discourse also speaks of separation and judgment. Chrysostom's confidence in the Gospel's power does not establish that every hearer will inevitably receive it or that every community will flourish without loss. The disciples must hear, inquire and act. His explanation of parables as an invitation to inquiry makes intelligent reception part of their task. The people who pass the word on are themselves still learning from Christ.\footnote{Matthew 13:36--50; Chrysostom, \textit{Homilies on Matthew}, 47.1.}

\begin{fourSenses}
\item[Literal.] Jesus describes growth and permeation; Paul gives thanks for a community whose reception, conduct and public example are visible amid affliction. The prayers ask God to form and sustain the people who worship him.
\item[Allegorical.] Christ gives the Gospel its power in the Church. Those who receive it are united in service of the living God, confessed in the Preface as Father, Son and Holy Spirit. The risen Son is both the one awaited and the source of deliverance.
\item[Moral.] Faith becomes work, charity becomes costly service and hope sustains endurance. Christians imitate their teachers insofar as they follow Christ, remain attentive to his word, forgive others and ask for the cleansing their witness requires.
\item[Anagogical.] The community awaits resurrection and final deliverance with the Son from heaven. The enduring praise of the Gradual and the desire expressed after Communion concern a completion beyond the present reach of its preaching.
\end{fourSenses}
